\documentclass[10pt,a4paper]{article}

\usepackage[utf8]{inputenc}
\usepackage[T1]{fontenc}
\usepackage[margin=0.65in]{geometry}
\usepackage{amsmath}
\usepackage{array}
\usepackage{booktabs}
\usepackage{enumitem}
\usepackage{longtable}
\usepackage{microtype}
\usepackage{tabularx}
\usepackage{xurl}
\usepackage[colorlinks=true,urlcolor=blue!55!black]{hyperref}

% Generated by paper/manuscript/generate_results.py; do not edit by hand.
% Canonical baseline source: release/v0.1.0/experiments/baselines/comparison/system_comparison.json
% Canonical baseline source SHA256: 06fdafcb4f2bb2f65e29d52093549c38c1c17a5ea8c1c12e3d84dae2b0b11395
% Values were cross-checked against every per-system scores/summary.json artifact.
\newif\ifllmresultsavailable
\llmresultsavailablefalse

\newcommand{\NumCards}{91}
\newcommand{\BudgetK}{10}
\newcommand{\BestQSARUtility}{74.966}
\newcommand{\BestQSARNDCG}{0.938}
\newcommand{\BestQSARRegret}{4.596}
\newcommand{\SimilarityUtility}{73.288}
\newcommand{\RandomUtility}{68.469}
\newcommand{\RulesUtility}{66.922}
\newcommand{\OracleUtility}{79.563}

% No canonical corrected LLM comparison artifact exists; keep the paper in pre-run mode.
\newcommand{\BestLLMLabel}{PENDING}
\newcommand{\BestLLMUtility}{PENDING}
\newcommand{\BestLLMNDCG}{PENDING}
\newcommand{\BestLLMActionValidity}{PENDING}
\newcommand{\BestLLMCost}{PENDING}

% Generated by paper/manuscript/generate_results.py; do not edit by hand.
% Canonical sources are the six matched raw/post-hoc-repaired trace pairs:
% release/v0.1.0/experiments/llm/matrix/openai_gpt_5_5_2026_04_23_selector/bare_llm/trace.jsonl SHA256: 90511fd8a98e693193eec6040141564d09061f3ca954464732c0ec0a84103b51
% release/v0.1.0/experiments/llm/matrix/openai_gpt_5_5_2026_04_23_selector/bare_llm/posthoc_repair.trace.jsonl SHA256: dbae88627f651c30250ac32166739ca995a6278705768295edb6e16b41b70609
% release/v0.1.0/experiments/llm/matrix/openai_gpt_5_5_2026_04_23_selector/llm_tools/trace.jsonl SHA256: a12c93417f0cabcad5d5af17c7d312ffd695f5f61dfe114225252551625f409e
% release/v0.1.0/experiments/llm/matrix/openai_gpt_5_5_2026_04_23_selector/llm_tools/posthoc_repair.trace.jsonl SHA256: 22c29705e1f2de103e14313b17d9be33f58ca8e58c198867139782b0a7fa9e21
% release/v0.1.0/experiments/llm/matrix/anthropic_opus_4_8_selector/bare_llm/trace.jsonl SHA256: a15731076cf24f16735025202c5e6524c9580eea0625aa24e89104d04ea1688f
% release/v0.1.0/experiments/llm/matrix/anthropic_opus_4_8_selector/bare_llm/posthoc_repair.trace.jsonl SHA256: 6e86077765501b2f335aa72107cae18508c51f8426f4e84e9174e8ce5fe83bc9
% release/v0.1.0/experiments/llm/matrix/anthropic_opus_4_8_selector/llm_tools/trace.jsonl SHA256: 4417c966d24ae9c599c351da45c52f8d11cc7a73fd7bd65f1965d2ec945130cb
% release/v0.1.0/experiments/llm/matrix/anthropic_opus_4_8_selector/llm_tools/posthoc_repair.trace.jsonl SHA256: 58dc763d1ca05a0b5728801a8c9ffd9abc4f3f1947bed85f5ffa0698f005a017
% release/v0.1.0/experiments/llm/matrix/deepseek_v4_pro_2026_07_16_selector/bare_llm/trace.jsonl SHA256: 299af81ecf1b511229278c1e696c73aff3bae05de66ecc4b84f3bc0e30a06829
% release/v0.1.0/experiments/llm/matrix/deepseek_v4_pro_2026_07_16_selector/bare_llm/posthoc_repair.trace.jsonl SHA256: 07a605ac55284724d6b93f0fbbf04d6437dd71e2cdb508afee26e6c623c11f61
% release/v0.1.0/experiments/llm/matrix/deepseek_v4_pro_2026_07_16_selector/llm_tools/trace.jsonl SHA256: ba727569fc0fcfecff685c0f3ffa9532147a236f732a4b56492c0dbbc5bfefea
% release/v0.1.0/experiments/llm/matrix/deepseek_v4_pro_2026_07_16_selector/llm_tools/posthoc_repair.trace.jsonl SHA256: 35f9512dff9426f26a6f7129ec67a14b66da02764f84239d84aa995372afea75
\newcommand{\RepairAttributionRows}{%
OpenAI, basic & 19/91 (20.88\%) & 10/91 & 14/910 (1.54\%) & 0.74; 1 (0--4) & 9 & 0 \\%
OpenAI, descriptors & 23/91 (25.27\%) & 10/91 & 15/910 (1.65\%) & 0.65; 0 (0--4) & 13 & 0 \\%
Anthropic, basic & 74/91 (81.32\%) & 71/91 & 252/910 (27.69\%) & 3.41; 2 (0--10) & 3 & 1 \\%
Anthropic, descriptors & 69/91 (75.82\%) & 68/91 & 229/910 (25.16\%) & 3.32; 3 (0--10) & 1 & 1 \\%
DeepSeek, basic & 83/91 (91.21\%) & 81/91 & 534/910 (58.68\%) & 6.43; 7 (0--10) & 2 & 30 \\%
DeepSeek, descriptors & 69/91 (75.82\%) & 69/91 & 355/910 (39.01\%) & 5.14; 4 (1--10) & 0 & 12 \\%
}
\newcommand{\RepairAttributionCompactRows}{%
OpenAI, basic & 19/91 (20.88\%) & 10/91 & 14/910 (1.54\%) & 0 \\%
OpenAI, descriptors & 23/91 (25.27\%) & 10/91 & 15/910 (1.65\%) & 0 \\%
Anthropic, basic & 74/91 (81.32\%) & 71/91 & 252/910 (27.69\%) & 1 \\%
Anthropic, descriptors & 69/91 (75.82\%) & 68/91 & 229/910 (25.16\%) & 1 \\%
DeepSeek, basic & 83/91 (91.21\%) & 81/91 & 534/910 (58.68\%) & 30 \\%
DeepSeek, descriptors & 69/91 (75.82\%) & 69/91 & 355/910 (39.01\%) & 12 \\%
}


\setlist{nosep,leftmargin=*}
\newcolumntype{Y}{>{\raggedright\arraybackslash}X}
\makeatletter
\renewcommand\section{\@startsection{section}{1}{\z@}%
  {-3.0ex \@plus -0.8ex \@minus -.2ex}%
  {1.5ex \@plus .2ex}%
  {\normalfont\Large\bfseries}}
\makeatother
\DeclareMathSizes{9}{10}{7}{5}

\title{\textbf{Supplementary Methods and Artifact Record}\\
SpecGuard-Chem 0.1.0}
\author{Daniel Hussey}
\date{Final university-feedback revision, 28 July 2026}

\begin{document}
\maketitle
\small
\setlist[itemize]{label=\(\bullet\)}

\section{Release status}

This supplement accompanies the results manuscript \emph{Before the Lab Acts:
Benchmarking Language Models for Constrained Compound Selection}. The corrected
data, deterministic baselines, and full 546-response provider matrix are
complete, committed, and audited. Model comparisons use six complete 91-card
pre-repair traces and six zero-call post-hoc-repaired views. No historical
response or score from the earlier 50-card draft is reported here. Those
artifacts are invalid because positional split references were previously
misread as dataframe labels; they remain only as immutable provenance.
Archival tagging and publication remain separate release-administration
actions. Chemical-diversity and multi-round extensions are deferred to a
separately versioned future analysis.

\section{Artifact visibility contract}

\begin{table}[h]
\centering
\small
\begin{tabularx}{\linewidth}{@{}p{0.27\linewidth}p{0.15\linewidth}Y@{}}
\toprule
\textbf{Field} & \textbf{Visibility} & \textbf{Role} \\
\midrule
Task, target, endpoint, source, activity semantics & System & Assay context and higher-is-better pChEMBL identity \\
Support IDs, structures, descriptors, pChEMBL activity & System & Project-local evidence observed before selection \\
Candidate IDs, structures, permitted descriptors & System & Finite action space \\
Constraints and budget & System & Executable action contract \\
Candidate activity & Evaluator only & Retrospective utility, ranking, and regret \\
Raw source paths and build-only metadata & Neither & Construction provenance, not model evidence \\
\bottomrule
\end{tabularx}
\end{table}

The public JSONL and scorer JSONL share benchmark, data, source, and transform
provenance. Every scorer row carries the SHA256 digest of the canonical public
card. The evaluator verifies identical task sets, candidate identity and order,
provenance, and binding hashes before hydration. Non-oracle systems receive a
strict allowlisted projection even when legacy monolithic fixtures are used.

\section{Final data identity}

\begin{table}[h]
\centering
\small
\begin{tabularx}{\linewidth}{@{}p{0.34\linewidth}Y@{}}
\toprule
\textbf{Artifact} & \textbf{SHA256} \\
\midrule
Normalized corrected records & \path{cec651b6e97f044bf82820c465b441763cafa18a34b12361a33e79ab30faf438} \\
Build configuration & \path{425a40d3b64ac8398c49dfc04508e1c66522754c4e07e75e69810bc2517d9a6a} \\
System-input cards & \path{c18e66c726bb26f8afc3ba8422b21ec327444560d92750421f0dc44a2f393d9e} \\
Scorer outcomes & \path{96b5d6060e3c75dda34d835fd166fd074ca5621c18924aa0ea2714acba173ff4} \\
Card-build metadata & \path{d986ba96589032c59dac2dcc24cadf9aa3325616fa2a395cec682ece7220af54} \\
Card-build audit & \path{abb80ad110ef71d69a2cea2f09cc456399a81d431bc9e365320ab8b27064812c} \\
Exact 546-request export & \path{50e518893b19d4a7efd64c62e08ab94d610815f8fb7518c9af4b64ff40b6f6c5} \\
Pre-run cost estimate & \path{d11ce35da68e5082153be4bc57c027915ccb76cf8fb19b02e5adf767b2ab525d} \\
\bottomrule
\end{tabularx}
\end{table}

An independent final rebuild produced byte-identical system input, scorer
outcomes, build metadata, and build audit. A second baseline run produced
byte-identical per-system traces and score directories. A second request export
produced the same 546 rows and request-file digest.

\section{Task inclusion and exclusions}

All 100 official \texttt{LO\_All} task keys were considered. The inclusion rule
keeps every task with at least ten candidates satisfying the declared candidate
constraints. Ninety-one tasks are included. The following tasks were excluded
without outcome-based selection:

\begin{longtable}{@{}lr@{}}
\toprule
\textbf{CARA task key} & \textbf{Feasible candidates} \\
\midrule
\endhead
CHEMBL1827734\_IC50 & 0 \\
CHEMBL3705332\_IC50 & 0 \\
CHEMBL3707531\_IC50 & 0 \\
CHEMBL3887078\_IC50 & 0 \\
CHEMBL3887456\_EC50 & 0 \\
CHEMBL3888518\_Ki & 0 \\
CHEMBL3888871\_Ki & 0 \\
CHEMBL3889212\_IC50 & 6 \\
CHEMBL3889266\_IC50 & 8 \\
\bottomrule
\end{longtable}

The included endpoint distribution is 68 IC50, 17 Ki, 5 EC50, and one Potency
task. Each included task has 50 support measurements and a distinct CARA target
identifier. No ranking by candidate-pool size or other purposive subsampling was
used. The feasibility rule is necessary for an exact top-10 action, but it
limits inference to these 91 CARA lead-optimization cards; it does not establish
representativeness of virtual-screening tasks, the excluded cards, or
prospective projects.

\section{Constraints}

Output validity requires exactly ten unique candidate IDs, membership in the
supplied candidate pool, and support-set exclusion. Candidate eligibility
requires an RDKit-parseable structure, molecular weight no greater than 500 Da,
and RDKit Crippen cLogP no greater than 4.5. The molecular-weight limit follows
the conventional rule-of-five boundary. The prespecified cLogP cutoff is a
simplified benchmark threshold, 0.5 below the conventional value of 5.0; it was
fixed before model evaluation. It is not a universally validated
medicinal-chemistry cutoff. The configured forbidden-SMARTS list is empty.

\begin{table}[h]
\centering
\small
\caption{Mutually exclusive candidate-constraint outcomes across all 19,588
official \texttt{LO\_All} candidate records.}
\label{tab:constraint_accounting}
\begin{tabular}{@{}lrr@{}}
\toprule
\textbf{Outcome} & \textbf{Candidates} & \textbf{Percentage} \\
\midrule
Passed both thresholds & 10,039 & 51.25\% \\
Molecular weight $>500$ only & 2,433 & 12.42\% \\
cLogP $>4.5$ only & 3,765 & 19.22\% \\
Exceeded both & 3,351 & 17.11\% \\
Unparseable structure & 0 & 0.00\% \\
\bottomrule
\end{tabular}
\end{table}

These rules are transparent eligibility controls for specification adherence,
not claims of synthesis feasibility, oral bioavailability, selectivity, ADMET
suitability, toxicity, or medicinal-chemistry approval.

\section{Systems and frozen provider matrix}

Deterministic systems comprise a random-valid selector, descriptor-desirability
rules, similarity to the best observed support compound, per-task random-forest,
gradient-boosting and linear support-vector regressors, and a non-deployable
hidden-outcome oracle. Similarity and QSAR use 2,048-bit radius-2 Morgan
fingerprints. Each regressor is fitted separately on one card's 50 visible
support fingerprints and pChEMBL measurements. Random forest uses 100 trees,
minimum leaf size 1, and seed 7. Gradient boosting uses the locked
scikit-learn defaults and seed 7. Linear SVR uses a linear kernel, $C=1.0$, and
\texttt{StandardScaler(with\_mean=False)}. There is no card-specific
hyperparameter search. Predictions rank all feasible candidates in descending
order; candidate ID breaks ties. Hidden candidate values are used only by the
scorer and oracle.

The rules-only baseline and deterministic repair fallback use the same
outcome-blind desirability score:
{\normalsize
\begin{align*}
D(i)={}&0.45\max\!\left(0,1-\frac{|\mathrm{MW}_i-350|}{350}\right)\\
 &+0.35\max\!\left(0,1-\frac{|\mathrm{cLogP}_i-2.5|}{4.5}\right)\\
 &+0.20\max\!\left(0,1-\frac{|\mathrm{TPSA}_i-75|}{150}\right).
\end{align*}
}
Feasible candidates are sorted by decreasing $D(i)$ with candidate ID as the
tie-break. RDKit-computed TPSA enters this software fallback even for the basic
LLM interface, which did not expose TPSA to the model. Random-valid shuffling
seeds Python's random-number generator with the integer represented by the first
eight hexadecimal digits of the SHA256 digest of seed 7, task ID, and
\texttt{random\_valid}, joined in that order by colons.

The provider matrix contains three conditions and two representations across 91
cards (546 requests):

\begin{itemize}
  \item OpenAI \texttt{gpt-5.5-2026-04-23}, low reasoning effort and provider
  JSON-object response mode;
  \item Anthropic \texttt{claude-opus-4-8}, no extended thinking, using an
  explicit JSON-only prompt and deterministic one-object extraction; and
  \item DeepSeek \texttt{deepseek-v4-pro}, thinking disabled and provider
  JSON-object response mode, as checked on 16 July 2026, with the
  provider-returned model identifier retained in traces.
\end{itemize}

The basic interface supplies IDs, structures, molecular weight, and cLogP. The
descriptor-enriched interface additionally supplies TPSA, hydrogen-bond donor
and acceptor counts, and rotatable-bond count. Both receive the same action
contract and an explicit instruction that pChEMBL is higher-is-better. The
full ordered candidate pool is supplied on every card. Candidate outcomes are
never supplied. Each condition uses a 4,096-token output limit; temperature and
provider seed are unset. Provider SDK retries are disabled. Each recorded
response therefore has one provider attempt, and malformed content remains
part of pre-repair behavior rather than triggering a paid retry. The
conservative uncached pre-run estimate is \$106.06, with maximum estimated
input 158,274 tokens and a hard cost ceiling of \$120.

The staged pilot completed all six calls with one recorded attempt each, full
provenance, and \$0.449700535 aggregate usage-derived cost; cache-only replay
reproduced every score artifact. The later authorized matrix completed all 546
responses. Each of the six raw conditions has one successful provider attempt
per card and complete model, response, usage, latency, request-hash, and retained
content evidence. The six repaired traces cover the same 91 cards and add zero
provider calls. Usage-derived token-pricing cost is \TotalLLMCost{} dollars for
the full matrix, with complete cost coverage.

The locked execution environment used Python 3.12.3, RDKit 2026.3.1,
scikit-learn 1.8.0, NumPy 2.4.4, pandas 3.0.2, OpenAI SDK 2.36.0, and
Anthropic SDK 0.100.0. The complete transitive environment and package hashes
are recorded in \texttt{uv.lock}; exact model requests and cached responses are
release artifacts.

\section{Deterministic post-hoc repair}

Repair is computed from the same raw response and adds zero provider calls:

\begin{enumerate}
  \item Parse and validate the raw output against the action contract.
  \item Retain valid raw candidate IDs in their original order.
  \item Remove duplicates, out-of-pool IDs, support identities, and infeasible
  candidates.
  \item Fill unoccupied slots using the deterministic rules-only ranking.
  \item Revalidate and store both outputs and issue sets, the source-trace
  digest, the repair policy, and a zero-provider-call attribution field.
\end{enumerate}

Raw metrics characterize model behavior. Repaired metrics characterize the
model-plus-harness system and must never be relabeled as unaided model ability.

``Cards repaired'' counts any contract modification. ``Fallback positions''
counts final candidate identities that were not retained from the valid,
unique pre-repair sequence and were instead supplied by the fixed rules
ranking. A card can therefore be repaired without receiving a fallback
identity.

\begin{table}[h]
\centering
\caption{Complete deterministic repair attribution. The per-repaired-card
column reports mean; median (range) fallback positions.}
\label{tab:supp_repair_attribution}
\scriptsize
\begin{tabularx}{\linewidth}{@{}Yrrrrrr@{}}
\toprule
\textbf{Condition} &
\textbf{Repaired} &
\textbf{Identity changed} &
\textbf{Fallback positions} &
\textbf{Per repaired card} &
\textbf{Zero fallback} &
\textbf{All 10 fallback} \\
\midrule
\RepairAttributionRows
\bottomrule
\end{tabularx}
\end{table}

\section{Utility and statistical analysis}

Let $s_{mcr}$ be the ID returned by system $m$ for card $c$ at position $r$.
Define $a_{mcr}$ as $y_{s_{mcr},c}$ only when that position exists, the ID is a
feasible candidate with an available hidden outcome, and the ID has not appeared
at an earlier scored position; otherwise $a_{mcr}=0$. Only the first $k$
positions are scored, and feasible utility is
{\normalsize
\[
U_m(c)=\sum_{r=1}^{k}a_{mcr},
\]
}
so only the first occurrence of an otherwise valid ID can score. Later
duplicates, out-of-pool or infeasible IDs, unavailable outcomes, missing
positions, and positions beyond $k$ receive no credit. All labels have the
common pChEMBL negative-log molar direction, but endpoint type and assay
conditions differ. Absolute utility is therefore not interpreted as a
biologically interchangeable cross-assay quantity.

System contrasts are paired within card. For systems $A$ and $B$, the analysis
first calculates $\Delta_c=U_A(c)-U_B(c)$ after both systems saw the same
support set and candidate pool, then averages the 91 equally weighted
differences. For valid fixed-size top-10 actions, a card-specific additive
offset cancels from this contrast. NDCG@10 supplies a secondary within-card
normalization against the ideal feasible ranking.

Marginal confidence intervals use 1,000 card-level bootstrap samples and
NumPy's default random-number generator with seed 7. Each sample draws 91 cards
with replacement, and the 2.5th and 97.5th percentiles of the sample means form
the 95\% interval. Headline contrasts use 2,000 paired bootstrap samples with
seed 13: the 91 observed card-level differences are resampled with replacement,
preserving pairing, and percentile limits are taken from the bootstrap mean
differences. The reported proportion of bootstrap means above zero is a
bootstrap replicate proportion, not a posterior probability. Intervals do not
correct task selection, dependence beyond the card, repeated-call variation,
or post-selection of the displayed best systems.

\end{document}
